\subsection{Formalization and Verification}\label{formalization-and-verification}

The Lean formalization certifies the paper's model, admissibility contract, and main theorem chain in the finite explicit setting. Mechanization therefore provides a formal check that the stated impossibility and optimality claims follow from the declared regime and that no unstated side channels are required. For transparency we keep the appendix concise and provide the full proof artifacts as a supplementary formalization release.

The detailed Lean module table is provided in supplementary material. It lists the 10 core modules with line counts, theorem counts, and purposes.

\paragraph{What is in scope in the mechanization.}
The formalization covers the abstract observer model, the information barrier, adaptive fiberwise side-information budgets, task-level sufficiency under fixed representations, helper-view fusion laws, exact product-fiber laws for factorized representations, structured axis families with derivability-aware closure, and the zero-error frontier claims stated in the main text. The structured-axis layer includes closure, exchange, equicardinality, and basis results. The converse chain, graph-confusability threshold and block-scaling results, adaptive side-information layer, and task-level laws are all represented in the cited theorem stack, including positive-distortion and learned-representation finite explicit corollaries. The Lean development works in the finite explicit setting used by the theorem statements: finite class sets, finite attribute families, and explicit profile maps or induced finite relations. It certifies the converse, block-scaling laws, adaptive fiberwise lower and upper bounds, task-level sufficiency, helper-view fusion laws, factorized product laws, the finite deterministic distortion floor, and the restricted structured-axis matroid bridge in that setting rather than the complexity of estimating $A_\pi$ from an implicit representation such as a circuit, program, or learned model. \leanmeta{\LH{L1}, \LHrng{L}{4}{5}, \LHrng{LWDC}{1}{3}, \LH{GPH13}, \LH{GPH16}, \LHrng{GPH}{18}{19}, \LH{GPH24}, \LHrng{GPH}{27}{28}, \LH{GPH30}, \LHrng{GPH}{32}{43}}

\paragraph{How the backing is exposed in the paper.}
The main text uses inline handle tags as claim-level links into the formalization, and the automatically generated handle index records the corresponding Lean declarations. The claims-only verification pass for the current release marks 10 main-text claims with explicit proof paths, so the emphasis is on which claims are backed rather than on aggregate artifact size. \leanmeta{\LH{GPH35}}

\paragraph{What is moved to supplementary artifact.}
Implementation-specific operational details and extended code listings are included in supplementary material and are not required to follow the IT contribution in the main paper.

\subsection{Attribute-Only Formalization}\label{interface-only-formalization}

Attribute-only observation is formalized by an equivalence relation on values induced by observable query responses.

\begin{lstlisting}[style=lean]
structure InterfaceValue where
  fields : List (String * Nat)
deriving DecidableEq

def getField (obj : InterfaceValue) (name : String) : Option Nat :=
  match obj.fields.find? (fun p => p.1 == name) with
  | some p => some p.2 | none => none

def interfaceEquivalent (a b : InterfaceValue) : Prop :=
  forall name, getField a name = getField b name

def InterfaceRespecting (f : InterfaceValue -> a) : Prop :=
  forall a b, interfaceEquivalent a b -> f a = f b
\end{lstlisting}

\subsection{Corollary 6.3: Provenance Impossibility}\label{corollary-6.3-interface-only-cannot-provide-provenance}

Under attribute-only observation, provenance is constant on attribute-equivalence classes; therefore provenance cannot be recovered when distinct classes collide under the observable profile.

\begin{lstlisting}[style=lean]
theorem interface_provenance_indistinguishable
    (getProvenance : InterfaceValue -> Option DuckProvenance)
    (h_interface : InterfaceRespecting getProvenance)
    (obj1 obj2 : InterfaceValue)
    (h_equiv : interfaceEquivalent obj1 obj2) :
    getProvenance obj1 = getProvenance obj2 :=
  h_interface obj1 obj2 h_equiv
\end{lstlisting}

This is the mechanized form of the main-text impossibility statement: if an observer factors through attribute profile alone, it cannot separate equal-profile values by source/provenance.

\subsection{Abstract Model Lean Formalization}\label{abstract-model-lean-formalization}

The abstract model is formalized directly at the axis level and then connected to concrete instantiations.

\begin{lstlisting}[style=lean]
-- Axis-indexed representation
abbrev Typ (A : Finset Axis) := (a : Axis) -> a \in A -> axisType a

-- Two-axis setting used in the paper
abbrev Typ2 := Typ ({Axis.Bases, Axis.Shape} : Finset Axis)

-- Projectors
abbrev projBases (t : Typ2) := t Axis.Bases (by simp)
abbrev projShape (t : Typ2) := t Axis.Shape (by simp)
\end{lstlisting}

The corresponding isomorphism theorem establishes that the two-axis representation is complete for in-scope observables in the formal model.

\subsection{Reproducibility}

The full Lean development is provided in supplementary material. To verify locally:
\begin{enumerate}
\item Install Lean 4 and Lake (\url{https://leanprover.github.io/}).
\item From the release package root, run:
\begin{lstlisting}[style=lean]
cd proofs
lake build
\end{lstlisting}
\item Confirm successful build with no \texttt{sorry} placeholders.
\end{enumerate}
